% Provenance (never printed): NSI premiere "Algorithmique: cout des algorithmes (lineaire, quadratique),
% recherche dichotomique: cout logarithmique (here only counted as halvings)". See sources/book2/BRIEFS.md.
\chapter{The Cost of an Algorithm}\label{ch:g11:cost-of-an-algorithm}

Two programs look for the names that appear twice in a list. They give the same answer. On a hundred names both
answer at once; on a million names, one answers in well under a second and the other takes hours. Nothing in the
answers tells them apart. What tells them apart is how much work each does, and how that work grows with the size of
the list.

The previous chapter proved that a loop stops and gives the right answer. This one counts how many steps it takes
before it stops: exactly, by counting operations in the code, and then roughly, by timing it. The two agree in a way
that lets one predict, from a small run, how long a large one will take.

\begin{recall}
Linear search, binary search, selection sort and insertion sort, with their invariants and variants, are in chapter
\ref{ch:g11:proving-algorithms-correct}. Counting the comparisons a search or a sort makes, on small lists and by
hand, was done in Book 1, grade 8.
\end{recall}

\section{Counting elementary operations}\label{sec:g11:cost-of-an-algorithm:counting}

\begin{definition}[Elementary operation, input size]\label{def:g11:cost-of-an-algorithm:elementary}
An \emph{elementary operation}\index{elementary operation} is an operation whose time does not depend on the size of
the data: comparing two numbers, adding them, reading or writing one element of a list. The \emph{input
size}\index{input size} is a number that measures the data, usually written $n$: the length of a list, the number of
names.
\end{definition}

\begin{definition}[Cost]\label{def:g11:cost-of-an-algorithm:cost}
The \emph{cost}\index{cost (time cost)} of an algorithm on an input is the number of \omterm{def:g11:cost-of-an-algorithm:elementary}{elementary operations} it
performs. One usually counts a single kind of operation that the others follow, such as the comparisons between
elements in a search or a sort.
\end{definition}

The simplest way to count is to let the program count: a variable that grows by one at each comparison.

\omcode{code/grade-11/06-cost-of-an-algorithm/listings/count_linear.py}{1}{18}{Linear search, counting its
comparisons.}\label{lst:g11:cost-of-an-algorithm:count-linear}

\omoutput{code/grade-11/06-cost-of-an-algorithm/outputs/count_linear.txt}{The first name \omterm{def:g11:cost-of-an-algorithm:cost}{costs} one comparison; a name
that is absent \omterm{def:g11:cost-of-an-algorithm:cost}{costs} $n$.}\label{out:g11:cost-of-an-algorithm:count-linear}

\section{Best and worst case}\label{sec:g11:cost-of-an-algorithm:cases}

\begin{definition}[Best case, worst case]\label{def:g11:cost-of-an-algorithm:cases}
For a given \omterm{def:g11:cost-of-an-algorithm:elementary}{input size} $n$, the \emph{best case}\index{best case} of an algorithm is an input of size $n$ on which its
\omterm{def:g11:cost-of-an-algorithm:cost}{cost} is the smallest, and the \emph{worst case}\index{worst case} one on which it is the largest. The \omterm{def:g11:cost-of-an-algorithm:cost}{cost} in the worst
case is a promise: no input of that size \omterm{def:g11:cost-of-an-algorithm:cost}{costs} more.
\end{definition}

Insertion sort shows how far apart the two may be. On a sorted list, each new element is compared once with the one
before it and stays put. On a list in reverse order, each new element is compared with every element before it and
travels to the front.

\omcode{code/grade-11/06-cost-of-an-algorithm/listings/count_insertion.py}{1}{26}{Insertion sort, counting the
comparisons between elements.}\label{lst:g11:cost-of-an-algorithm:count-insertion}

\omoutput{code/grade-11/06-cost-of-an-algorithm/outputs/count_insertion.txt}{A sorted list \omterm{def:g11:cost-of-an-algorithm:cost}{costs} $n - 1$ comparisons,
a reversed one $n(n-1)/2$, and a shuffled one about half of that.}\label{out:g11:cost-of-an-algorithm:count-insertion}

\begin{omfigure}
\begin{tikzpicture}[c/.style={draw, minimum size=7mm, font=\small}]
  \node[font=\footnotesize, anchor=east] at (-0.6,0) {sorted};
  \node[font=\footnotesize, anchor=east] at (-0.6,-1.6) {reversed};
  \foreach \v [count=\k from 0] in {1, 2, 3, 4, 5, 6} \node[c, fill=omMeth!12] (a\k) at (0.7*\k, 0) {\v};
  \foreach \v [count=\k from 0] in {6, 5, 4, 3, 2, 1} \node[c, fill=omThm!12] (b\k) at (0.7*\k, -1.6) {\v};
  \foreach \k in {1,...,5} {
    \pgfmathtruncatemacro{\p}{\k - 1}
    \draw[->, omDef, thick] (a\k.north) to[bend right=50] (a\p.north);
  }
  \draw[->, omThm, thick] (b5.south) to[bend left=25] node[below, font=\footnotesize] {$5$ comparisons} (b0.south);
  \node[font=\footnotesize, anchor=west] at (4.4,0) {each new element: $1$ comparison, $5$ in all};
  \node[font=\footnotesize, anchor=west] at (4.4,-1.6) {element $k$: $k$ comparisons, $1 + 2 + \dots + 5 = 15$};
\end{tikzpicture}

{\small Insertion sort on six elements. Sorted: every new element is checked once against its left neighbour. Reversed:
the last one, $1$, is compared with all five before it on its way to the front.}
\end{omfigure}

\section{Constant, linear and quadratic cost}\label{sec:g11:cost-of-an-algorithm:growth}

The exact count depends on details: counting assignments too, or starting the loop at $0$ or $1$, changes it a little.
What does not change is how it grows with $n$.

\begin{definition}[Constant, linear, quadratic cost]\label{def:g11:cost-of-an-algorithm:growth}
A \omterm{def:g11:cost-of-an-algorithm:cost}{cost} is \emph{constant}\index{constant, linear, quadratic cost} if it does not grow with $n$; \emph{linear} if it
is about proportional to $n$, so that doubling $n$ about doubles it; \emph{quadratic} if it is about proportional to
$n^2$, so that doubling $n$ about multiplies it by $4$.
\end{definition}

Linear search \omterm{def:g11:cost-of-an-algorithm:cost}{costs} $n$ in its \omterm{def:g11:cost-of-an-algorithm:cases}{worst case}: linear. Insertion sort \omterm{def:g11:cost-of-an-algorithm:cost}{costs} $n(n-1)/2 = (n^2 - n)/2$: for large $n$ the
term $n^2/2$ is almost all of it, and the \omterm{def:g11:cost-of-an-algorithm:cost}{cost} is quadratic.

\omcode{code/grade-11/06-cost-of-an-algorithm/listings/doubling.py}{1}{14}{The worst-case counts when $n$
doubles.}\label{lst:g11:cost-of-an-algorithm:doubling}

\omoutput{code/grade-11/06-cost-of-an-algorithm/outputs/doubling.txt}{Doubling $n$ doubles the linear count and
multiplies the quadratic one by $4$.}\label{out:g11:cost-of-an-algorithm:doubling}

\begin{omfigure}
\begin{tikzpicture}
\begin{axis}[omrecord, width=11cm, height=6cm, xlabel={$n$}, ylabel={comparisons, worst case}, xmin=0, xmax=100,
    ymin=0, ymax=5000, legend pos=north west, scaled y ticks=false,
    yticklabel style={/pgf/number format/1000 sep={\,}}]
  \addplot[very thick, omDef] table[x=n, y=linear]
    {code/grade-11/06-cost-of-an-algorithm/figures/counts.dat};
  \addplot[very thick, omThm] table[x=n, y=quadratic]
    {code/grade-11/06-cost-of-an-algorithm/figures/counts.dat};
  \legend{linear search, insertion sort}
\end{axis}
\end{tikzpicture}

{\small Worst-case counts, from \code{figures/counts.py}. At $n = 100$, $100$ comparisons against $4950$; the
straight line is not zero, it only looks flat beside the parabola.}
\end{omfigure}

\section{The cost of halving}\label{sec:g11:cost-of-an-algorithm:halving}

Binary search does not look at the elements one by one: it halves the window each time. Its \omterm{def:g11:cost-of-an-algorithm:cost}{cost} is the \omterm{def:g11:cost-of-an-algorithm:halvings}{number of
halvings}.

\begin{definition}[Number of halvings]\label{def:g11:cost-of-an-algorithm:halvings}
For a whole number $n \ge 1$, the \emph{number of halvings}\index{number of halvings h(n)} $h(n)$ is the number of
times $n$ can be halved, rounding down, before it reaches $1$. It is the whole number $h$ such that
$2^h \le n < 2^{h + 1}$.
\end{definition}

\omcode{code/grade-11/06-cost-of-an-algorithm/listings/halvings.py}{1}{14}{$h(n)$, counted, and checked against the
powers of $2$.}\label{lst:g11:cost-of-an-algorithm:halvings}

\omoutput{code/grade-11/06-cost-of-an-algorithm/outputs/halvings.txt}{A thousand is halved $9$ times, a million $19$,
a billion $29$: multiplying $n$ by a thousand adds about $10$.}\label{out:g11:cost-of-an-algorithm:halvings}

\begin{omfigure}
\begin{tikzpicture}
\begin{axis}[omrecord, ybar, width=11cm, height=5.2cm, bar width=7mm, xlabel={halving}, ylabel={$n$}, ymin=0,
    ymax=1300, enlarge x limits=0.07, xtick={0,...,9}, nodes near coords, nodes near coords style={font=\scriptsize,
    /pgf/number format/1000 sep={}}]
  \addplot[fill=omDef!40, draw=omDef] coordinates {(0,1000) (1,500) (2,250) (3,125) (4,62) (5,31) (6,15) (7,7) (8,3)
    (9,1)};
\end{axis}
\end{tikzpicture}

{\small Halving $1000$, rounding down: after $9$ halvings only $1$ is left, so $h(1000) = 9$, and indeed
$2^9 = 512 \le 1000 < 1024$.}
\end{omfigure}

The binary search of the previous chapter looks at one element, \code{t[mid]}, per turn, comparing it with \code{x}
at most twice. Its window starts with $n$ elements, and each turn leaves at most half of it, rounded down. In the \omterm{def:g11:cost-of-an-algorithm:cases}{worst
case} it makes $h(n) + 1$ turns: $20$ for a million names, against a million comparisons for linear search. Its \omterm{def:g11:cost-of-an-algorithm:cost}{cost}
grows by one turn when $n$ doubles.

\section{Measured running times}\label{sec:g11:cost-of-an-algorithm:measured}

Counts predict running times, up to a factor that depends on the machine. \code{time.perf_counter()} reads a clock
meant for measuring short times; the difference between two readings is the time spent between them. One run is
never enough: the machine does other things at the same time, so each measurement is repeated and the median kept, the
middle value, which an occasional slow run does not move.

\omcode{code/grade-11/06-cost-of-an-algorithm/listings/measure_sorts.py}{1}{25}{Timing insertion sort as $n$
doubles.}\label{lst:g11:cost-of-an-algorithm:measure-sorts}

\omoutput{code/grade-11/06-cost-of-an-algorithm/outputs/measure_sorts.txt}{Measured on the author's laptop: each
doubling multiplies the time by about $4$, as the quadratic count
predicted.}\label{out:g11:cost-of-an-algorithm:measure-sorts}

\begin{omfigure}
\begin{tikzpicture}
\begin{axis}[omrecord, width=11cm, height=6cm, xlabel={$n$}, ylabel={milliseconds}, xmin=0, xmax=4200, ymin=0, ymax=210,
    scaled x ticks=false, xticklabel style={/pgf/number format/1000 sep={\,}}, xtick={0,1000,2000,3000,4000}]
  \addplot[very thick, omThm, mark=*, mark size=1.5pt] table[x=n, y=ms]
    {code/grade-11/06-cost-of-an-algorithm/figures/measured_plot.dat};
  \draw[dashed, omDef] (axis cs:1000,0) -- (axis cs:1000,11) (axis cs:2000,0) -- (axis cs:2000,45)
    (axis cs:4000,0) -- (axis cs:4000,187);
\end{axis}
\end{tikzpicture}

{\small Insertion sort on shuffled lists, median of five runs, measured on the author's laptop
(\code{figures/measured_plot.py}). The dashed lines mark $n = 1000$, $2000$, $4000$: each doubling, about four times
the time.}
\end{omfigure}

The same machine, looking for a number in a million: a scan of the list against a binary search and against a set,
which finds an element by computing where it is stored rather than by comparing.

\omcode{code/grade-11/06-cost-of-an-algorithm/listings/measure_search.py}{1}{31}{One lookup in a million numbers, three
ways.}\label{lst:g11:cost-of-an-algorithm:measure-search}

\omoutput{code/grade-11/06-cost-of-an-algorithm/outputs/measure_search.txt}{Measured on the author's laptop. The scan is
tens of thousands of times slower than binary search on a million elements.}\label{out:g11:cost-of-an-algorithm:measure-search}

\begin{remark}
A timing measures everything between the two clock readings. A \code{print} inside the timed part, or building the
list there, is measured too. Timings also vary from run to run and from machine to machine: this book's timings are
printed as measured once, on one laptop, and only their ratios are meant to carry over.
\end{remark}

\section{Exercises}\label{sec:g11:cost-of-an-algorithm:exercises}

\begin{exercise}[$\star$]\label{exo:g11:cost-of-an-algorithm:1}
How many comparisons does linear search make, in its \omterm{def:g11:cost-of-an-algorithm:cases}{worst case}, on a list of $500$ names? When is that case?
\end{exercise}

\begin{exercise}[$\star$]\label{exo:g11:cost-of-an-algorithm:2}
Compute $h(1000)$ by halving by hand, and check it with powers of $2$.
\end{exercise}

\begin{exercise}[$\star$]\label{exo:g11:cost-of-an-algorithm:3}
What does \code{count_insertion.py} print for the sorted and the reversed lists if $n = 20$?
\end{exercise}

\begin{exercise}[$\star$]\label{exo:g11:cost-of-an-algorithm:4}
How many times does the body run in \code{for i in range(n)}, then inside it \code{for j in range(i)}? Give the count
for $n = 10$ and in general.
\end{exercise}

\begin{exercise}[$\star$]\label{exo:g11:cost-of-an-algorithm:5}
Write \code{count_comparisons(sort, t)} that sorts a copy of \code{t} with a given sorting function, such as
\code{sorted}, and returns how many comparisons it made. Hint: \code{functools.cmp_to_key} turns a comparison function
into a key.
\end{exercise}

\begin{exercise}[$\star\star$]\label{exo:g11:cost-of-an-algorithm:6}
Selection sort compares, at turn $i$, the element in place $i$ and those after it. Show that it makes exactly
$n(n-1)/2$ comparisons on every list, by writing the sum $(n-1) + (n-2) + \dots + 1$ twice, the second time backwards,
and adding.
\end{exercise}

\begin{exercise}[$\star\star$]\label{exo:g11:cost-of-an-algorithm:7}
A student times a sort with a \code{print} of the list after each turn, inside the timed part. Why is the measurement
wrong, and what does it measure?
\end{exercise}

\begin{exercise}[$\star\star$]\label{exo:g11:cost-of-an-algorithm:8}
A quadratic \omterm{def:g11:cost-of-an-algorithm:cost}{cost} takes $2$ seconds for $n = 10\,000$. About how long for $n = 30\,000$?
\end{exercise}

\begin{exercise}[$\star\star$]\label{exo:g11:cost-of-an-algorithm:9}
Write \code{has_duplicate(t)} that compares every pair of elements once, and give its number of comparisons when there
is no duplicate.
\end{exercise}

\begin{exercise}[$\star\star$]\label{exo:g11:cost-of-an-algorithm:10}
Which list of length $n$ is the \omterm{def:g11:cost-of-an-algorithm:cases}{worst case} for the chapter's \code{largest}, if one counts the assignments
\code{best = t[i]}? How many are there then?
\end{exercise}

\begin{exercise}[$\star\star$]\label{exo:g11:cost-of-an-algorithm:11}
Write \code{h(n)} by halving, and check on all $n$ from $1$ to $5000$ that $2^{h(n)} \le n < 2^{h(n) + 1}$.
\end{exercise}

\begin{exercise}[$\star\star\star$]\label{exo:g11:cost-of-an-algorithm:12}
Prove that the binary search of the previous chapter makes at most $h(n) + 1$ turns on a list of $n \ge 1$
elements.
\end{exercise}

\begin{exercise}[$\star\star\star$]\label{exo:g11:cost-of-an-algorithm:13}
Rank the \omterm{def:g11:cost-of-an-algorithm:cost}{costs} $1000$, $h(n)$, $100n$, $n \cdot h(n)$ and $n^2$ at $n = 10$, $1000$ and $10^6$. Which is smallest
at $n = 10$? Is $n \cdot h(n)$ below $100n$ at $n = 10^6$? What does this say about comparing \omterm{def:g11:cost-of-an-algorithm:cost}{costs} only by their
growth?
\end{exercise}

\begin{exercise}[$\star\star\star$]\label{exo:g11:cost-of-an-algorithm:14}
Write a duplicate search that sorts the list and then compares neighbours. How many comparisons does the scan make?
\end{exercise}

\begin{exercise}[$\star\star\star$]\label{exo:g11:cost-of-an-algorithm:15}
From the three times of \code{measure_sorts.py}, predict the time of insertion sort for $n = 40\,000$. Running it
would take a while: estimate how long, and say why the prediction might be off.
\end{exercise}

\section{Problem: Two Ways to Find Duplicates}\label{sec:g11:cost-of-an-algorithm:problem}

\begin{problem}[{Weekend problem --- from which size sorting first pays}]\label{pb:g11:cost-of-an-algorithm:1}
A list of $n$ numbers may hold a duplicate. Three functions answer: one compares every pair, one sorts the list with
Python's \code{sorted} and compares neighbours, one builds a set. The first two are counted with a value that counts
each comparison made with it; the third compares nothing, it computes where each number goes.

\omcode{code/grade-11/06-cost-of-an-algorithm/listings/duplicates.py}{1}{34}{Three ways to find a duplicate, and a
counted value.}\label{lst:g11:cost-of-an-algorithm:duplicates}

\omcode{code/grade-11/06-cost-of-an-algorithm/listings/dup_count.py}{1}{27}{Counting the comparisons on lists without
duplicates.}\label{lst:g11:cost-of-an-algorithm:dup-count}

\omoutput{code/grade-11/06-cost-of-an-algorithm/outputs/dup_count.txt}{The counts, on shuffled lists of different
numbers.}\label{out:g11:cost-of-an-algorithm:dup-count}

\omcode{code/grade-11/06-cost-of-an-algorithm/listings/dup_measure.py}{1}{30}{Timing the three, and predicting the
quadratic one at $n = 100\,000$.}\label{lst:g11:cost-of-an-algorithm:dup-measure}

\omoutput{code/grade-11/06-cost-of-an-algorithm/outputs/dup_measure.txt}{Measured on the author's
laptop.}\label{out:g11:cost-of-an-algorithm:dup-measure}

\needspace{4\baselineskip}\medskip\noindent\textbf{Part I --- Every pair.}
\begin{enumerate}
  \item Why is a list without duplicate the \omterm{def:g11:cost-of-an-algorithm:cases}{worst case} of \code{by_pairs}?
  \item Count its comparisons on such a list of $n$ elements, and check the line $n = 100$.
  \item Is its \omterm{def:g11:cost-of-an-algorithm:cost}{cost} \omterm{def:g11:cost-of-an-algorithm:growth}{constant}, linear or quadratic?
  \item What does \code{Counted} do, and why can it count the comparisons made inside \code{sorted}?
  \item Why is \code{by_pairs} not run at $n = 100\,000$ in \code{dup_count.py}?
\end{enumerate}

\needspace{4\baselineskip}\medskip\noindent\textbf{Part II --- Sort, then scan.}
\begin{enumerate}[resume]
  \item Why do two equal values end up side by side once the list is sorted?
  \item How many comparisons does the scan make, at most?
  \item Subtract them from the printed counts at $n = 10$ and $n = 100$: how many comparisons did the sort make?
  \item From the output, from which $n$ does sorting win, on these lists? Why did it lose for small $n$?
  \item At $n = 100\,000$, how many times fewer comparisons than every pair?
\end{enumerate}

\needspace{4\baselineskip}\medskip\noindent\textbf{Part III --- The measured times.}
\begin{enumerate}[resume]
  \item Check that the times of \code{by_pairs} at $1000$, $2000$, $4000$ grow like a quadratic \omterm{def:g11:cost-of-an-algorithm:cost}{cost}.
  \item Explain the factor $625$ in the prediction for $n = 100\,000$.
  \item Express the predicted time in minutes.
  \item The set version compares nothing. Why can it still not take a \omterm{def:g11:cost-of-an-algorithm:growth}{constant} time?
  \item Rank the three by measured time at $n = 100\,000$.
\end{enumerate}

\needspace{4\baselineskip}\medskip\noindent\textbf{Part IV --- The verdict.}
\begin{enumerate}[resume]
  \item Check the printed speed-up from the two times it divides.
  \item Why is the speed-up measured in time different from the ratio of comparisons of question 10?
  \item For a list of a thousand numbers checked once, does the choice matter? For a million?
  \item Why is the prediction for \code{by_pairs} at $100\,000$ reasonable but not exact?
  \item State the result: the size from which sorting needs fewer comparisons, and the speed-up measured at
    $n = 100\,000$.
\end{enumerate}
\end{problem}
